% =========================================================================
% Capítulo 5: Mecánica Lagrangiana
% Tema 01: Mecánica Lagrangiana
% Problema 06: Péndulo doble de dos varillas rígidas
% =========================================================================
\subsubsection{Péndulo doble de dos varillas rígidas}

Un péndulo doble oscila libremente en un plano vertical bajo un campo gravitatorio
constante $g$. Está formado por dos barras delgadas y homogéneas de masas $m_1,m_2$ y
longitudes $\ell_1,\ell_2$, respectivamente, como se muestra en la Figura
\ref{fig:c5_e01_p06_pendulo_doble}. La primera barra está articulada por uno de sus
extremos a un pivote fijo sin fricción en el origen $O$; la segunda barra se acopla de
forma libre al extremo inferior de la primera. El movimiento se describe con las
coordenadas generalizadas $\{\theta_1,\theta_2\}$, medidas respecto a la vertical
descendente. Considere los momentos de inercia baricéntricos $I_i=\tfrac{1}{12}m_i\ell_i^2$.

\begin{enumerate}
    \item Identifique los grados de libertad y las coordenadas generalizadas.
    \item Determine las posiciones y velocidades de los centros de masa de ambas barras.
    \item Construya el Lagrangiano $\mathcal{L}=T-V$ del sistema.
    \item Deduzca las dos ecuaciones de movimiento acopladas usando las ecuaciones de
          Euler-Lagrange.
\end{enumerate}

\begin{figure}[!ht]
    \centering
    \begin{tikzpicture}[>=Latex, thick, scale=1.2]

        % --- SOPORTE FIJO SUPERIOR ---
        \fill[gray!40] (-0.6,0) rectangle (0.6,0.15);
        \draw (-0.6,0) -- (0.6,0);

        % --- COORDENADAS PRINCIPALES ---
        \coordinate (O) at (0,0);                            % Pivote fijo
        \coordinate (V1) at (0,-1.5);                        % Referencia vertical 1
        \coordinate (A) at (270-60:2.5);                     % Extremo barra 1
        \coordinate (G1) at (270-60:1.25);                   % Centro de masa barra 1
        \coordinate (V2) at ($(A) + (0,-1.5)$);              % Referencia vertical 2
        \coordinate (B) at ($(A) + (270-30:2.2)$);           % Extremo barra 2
        \coordinate (G2) at ($(A) + (270-30:1.1)$);          % Centro de masa barra 2

        % --- LINEAS DE REFERENCIA VERTICALES ---
        \draw[dashed, gray] (O) -- (0,-3.2);
        \draw[dashed, gray] (A) -- ($(A) + (0,-2.2)$);

        % --- BARRAS DELGADAS (CUERPO RIGIDO) ---
        \draw[line width=3.5pt, line cap=round, gris1!70!black] (O) -- (A);
        \draw[line width=3.5pt, line cap=round, gris1!70!black] (A) -- (B);

        % --- ARTICULACIONES ---
        \fill[black] (O) circle (0.08) node[above left] {$O$};
        \fill[black] (A) circle (0.08);

        % --- CENTROS DE MASA ---
        \fill[myred] (G1) circle (0.07) node[right=2pt, black] {$G_1\,(m_1,\ell_1,I_1)$};
        \fill[myred] (G2) circle (0.07) node[left=2pt, black] {$G_2\,(m_2,\ell_2,I_2)$};

        % --- ANGULOS ---
        \draw[->] (O)+(270:0.9) arc (270:210:0.9);
        \node at ($(O)+(238:1.2)$) {$\theta_1$};
        \draw[->] (A)+(270:0.9) arc (270:240:0.9);
        \node at ($(A)+(252:1.25)$) {$\theta_2$};

    \end{tikzpicture}
    \caption{Péndulo doble de dos barras homogéneas articuladas en serie.}
    \label{fig:c5_e01_p06_pendulo_doble}
\end{figure}

% --- SOLUCIÓN ---------------------------------------------------------------
\begin{solucion}{Péndulo doble de dos varillas rígidas}
\subsubsection*{Pregunta 1: Grados de libertad}

El sistema tiene dos cuerpos rígidos que se mueven en el plano vertical, articulados en
serie. Cada articulación introduce un grado de libertad rotacional, de modo que el sistema
posee \textbf{dos grados de libertad}, descritos por las coordenadas generalizadas
$\{\theta_1,\theta_2\}$ medidas respecto a la vertical descendente.

\subsubsection*{Pregunta 2: Posiciones y velocidades}

Con origen en $O$, eje $x$ horizontal y eje $y$ vertical hacia arriba, el extremo inferior
de la primera barra (punto $A$) y los centros de masa $G_1,G_2$ son:
$$\vec{r}_{G_1} = \tfrac{\ell_1}{2}\big(\sin\theta_1,\,-\cos\theta_1\big),\qquad
\vec{r}_{A} = \ell_1\big(\sin\theta_1,\,-\cos\theta_1\big),$$
$$\vec{r}_{G_2} = \Big(\ell_1\sin\theta_1 + \tfrac{\ell_2}{2}\sin\theta_2,\;
-\ell_1\cos\theta_1 - \tfrac{\ell_2}{2}\cos\theta_2\Big).$$

Derivando respecto al tiempo se obtienen los cuadrados de las rapideces de los centros de
masa:
$$v_{G_1}^2 = \left(\tfrac{\ell_1}{2}\right)^2\dot\theta_1^2,\qquad
v_{G_2}^2 = \ell_1^2\dot\theta_1^2 + \tfrac{\ell_2^2}{4}\dot\theta_2^2
+ \ell_1\ell_2\cos(\theta_1-\theta_2)\,\dot\theta_1\dot\theta_2.$$
Las velocidades angulares de las barras son $\omega_1=\dot\theta_1$ y $\omega_2=\dot\theta_2$.

\subsubsection*{Pregunta 3: Lagrangiano}

La energía cinética de cada barra combina traslación de su centro de masa y rotación
baricéntrica ($I_i=\tfrac{1}{12}m_i\ell_i^2$):
$$T_1 = \tfrac12 m_1 v_{G_1}^2 + \tfrac12 I_1\dot\theta_1^2 = \tfrac16 m_1\ell_1^2\dot\theta_1^2,$$
$$T_2 = \tfrac12 m_2 v_{G_2}^2 + \tfrac12 I_2\dot\theta_2^2
= \tfrac12 m_2\ell_1^2\dot\theta_1^2 + \tfrac16 m_2\ell_2^2\dot\theta_2^2
+ \tfrac12 m_2\ell_1\ell_2\cos(\theta_1-\theta_2)\,\dot\theta_1\dot\theta_2.$$
Sumando y agrupando:
$$T = \tfrac12\left(\tfrac13 m_1 + m_2\right)\ell_1^2\,\dot\theta_1^2
+ \tfrac12\cdot\tfrac13 m_2\ell_2^2\,\dot\theta_2^2
+ \tfrac12 m_2\ell_1\ell_2\cos(\theta_1-\theta_2)\,\dot\theta_1\dot\theta_2.$$

La energía potencial gravitatoria (nivel de referencia en $O$) es
$$V = -\left(\tfrac12 m_1 + m_2\right)g\ell_1\cos\theta_1 - \tfrac12 m_2 g\ell_2\cos\theta_2.$$

Por lo tanto, el Lagrangiano $\mathcal{L}=T-V$ es
$$\boxed{\begin{aligned}
\mathcal{L} ={}& \tfrac12\left(\tfrac13 m_1 + m_2\right)\ell_1^2\dot\theta_1^2
+ \tfrac16 m_2\ell_2^2\dot\theta_2^2
+ \tfrac12 m_2\ell_1\ell_2\cos(\theta_1-\theta_2)\dot\theta_1\dot\theta_2 \\
&+ \left(\tfrac12 m_1 + m_2\right)g\ell_1\cos\theta_1 + \tfrac12 m_2 g\ell_2\cos\theta_2
\end{aligned}}$$

\subsubsection*{Pregunta 4: Ecuaciones de movimiento}

Aplicando $\dfrac{d}{dt}\!\left(\dfrac{\partial\mathcal{L}}{\partial\dot\theta_i}\right)
-\dfrac{\partial\mathcal{L}}{\partial\theta_i}=0$ y simplificando los términos de
Coriolis, se obtienen las dos ecuaciones acopladas:
$$\boxed{\begin{aligned}
&\left(\tfrac13 m_1 + m_2\right)\ell_1^2\,\ddot\theta_1
+ \tfrac12 m_2\ell_1\ell_2\cos(\theta_1-\theta_2)\,\ddot\theta_2 \\
&\quad + \tfrac12 m_2\ell_1\ell_2\sin(\theta_1-\theta_2)\,\dot\theta_2^2
+ \left(\tfrac12 m_1 + m_2\right)g\ell_1\sin\theta_1 = 0
\end{aligned}}$$
$$\boxed{\begin{aligned}
&\tfrac13 m_2\ell_2^2\,\ddot\theta_2
+ \tfrac12 m_2\ell_1\ell_2\cos(\theta_1-\theta_2)\,\ddot\theta_1 \\
&\quad - \tfrac12 m_2\ell_1\ell_2\sin(\theta_1-\theta_2)\,\dot\theta_1^2
+ \tfrac12 m_2 g\ell_2\sin\theta_2 = 0
\end{aligned}}$$

Para oscilaciones pequeñas ($\theta_i\ll1$, $\cos\approx1$, $\sin\approx\theta$ y términos
cuadráticos en velocidad despreciables) el sistema se linealiza a la forma
$\mathbf{M}\ddot{\boldsymbol\theta}+\mathbf{K}\boldsymbol\theta=0$, cuyos autovalores
proporcionan los modos normales de vibración del péndulo doble.
\end{solucion}
